% ======================================================================
% Section 2 — The pair-sum equality theorem
% ======================================================================

\section{The pair-sum equality theorem}\label{sec:equality}

\subsection{Setup}

Fix $n\ge 2$ and speeds $V=(v_1,\dots,v_n)$, not necessarily distinct in
this section (no argument below uses distinctness; the conjecture adds it).
Normalize $\gcd(V)=1$. For $t\in\T$ write
\[
f_p(t)=\wnc{v_p t},\qquad
g(t)=\min_{1\le p\le n} f_p(t),\qquad
M(V)=\max_{t\in\T} g(t).
\]
The function $g$ is continuous on the compact circle, so the maximum is
attained, and $M(V)>0$: at $t=1/(2v_{\max})$ every $f_p$ equals
$v_p/(2v_{\max})\le \tfrac12$, whence $g\ge v_{\min}/(2v_{\max})>0$.

\begin{definition}[pairs, lattices, reduced optima]\label{def:pair}
For an unordered pair of indices $\{u,w\}$, $u\ne w$, the \emph{pair sum}
is $N=N_{uw}:=v_u+v_w\ (\ge 3)$, and the \emph{pair-sum lattice} of
$\{u,w\}$ is the set of times $\{k/N : 1\le k\le N-1\}$. The
\emph{reduced optimum} of the pair is
\[
\tau(u,w)\;:=\;\max_{1\le k\le N-1}\;
\min_{p\ne w}\;\Bigl\lVert \frac{v_p k}{N}\Bigr\rVert .
\]
This is the $(n{-}1)$-runner optimal loneliness of the effective residues
$(v_p\bmod N)_{p\ne w}$ on the cyclic group $\Z_N$: runner $w$ has been
deleted, and by Lemma~\ref{lem:mirror} below the deleted runner is carried
by its partner $u$ at every lattice time.
\end{definition}

\begin{theorem}[pair-sum equality; lossless reduction]\label{thm:equality}
For every $V$ as above,
\[
M(V)\;=\;\max_{\{u,w\}}\;\tau(u,w),
\]
the maximum over the $\binom{n}{2}$ unordered pairs. Moreover the maximum
defining $M(V)$ is attained at a time $t^\ast=k^\ast/N$ on the pair-sum
lattice of some pair $\{u,w\}$ with $1\le k^\ast\le N-1$; and if
$M(V)<\tfrac12$ the binding runners at $t^\ast$ can be chosen two-sided,
one at residue $m^\ast$ and one at residue $-m^\ast$ on $\Z_N$.
\end{theorem}

\begin{corollary}[lossless reformulation]\label{cor:tau}
$\mathrm{LRC}$-$_n$ is equivalent to the finite statement
$(\mathrm{TAU}$-$_n)$: for every $\gcd$-normalized speed set $V$ there is
a pair $\{u,w\}$ with
\[
(n+1)\;\cdot\;\max_{1\le k\le N-1}\min_{p\ne w}
\Bigl\lVert \frac{v_p k}{N}\Bigr\rVert \;\ge\; N=v_u+v_w .
\]
\end{corollary}

The corollary is immediate: $\mathrm{LRC}$-$_n$ asserts $M(V)\ge
1/(n{+}1)$, i.e.\ $\max_{\{u,w\}}\tau(u,w)\ge 1/T$ with $T=n{+}1$, which
unfolds to exactly $(\mathrm{TAU}$-$_n)$; the equivalence is an
\emph{identity} between a continuous maximization and a finite family of
finite maximizations, so no information is lost in either direction. The
proof of the theorem occupies the rest of the section.

\subsection{The mirror identity}

\begin{lemma}[mirror / pair collapse]\label{lem:mirror}
Let $N=v_u+v_w$. For every integer $k$:
\begin{enumerate}[leftmargin=1.8em,itemsep=1pt]
\item[(i)] $v_u k\equiv -\,v_w k \pmod N$, hence
  $\lVert v_u k/N\rVert=\lVert v_w k/N\rVert$;
\item[(ii)] $\min_{p}\lVert v_p k/N\rVert
  =\min_{p\ne w}\lVert v_p k/N\rVert$.
\end{enumerate}
\end{lemma}

\begin{proof}
(i) The two speeds sum to $N$, so
$v_u k+v_w k=Nk\equiv 0 \pmod N$; the two positions are reflections of
each other through the origin, and distances to the origin coincide.
(ii) The left side is a minimum over a superset of the index set on the
right, so it is at most the right side. For the reverse inequality,
split the full minimum as
$\min\bigl(\min_{p\notin\{u,w\}} f_p,\,f_u,\,f_w\bigr)$ and replace the
last two terms by $f_u$ alone, which equals them by (i); the resulting
expression is exactly the right side, since
$\{p\ne w\}=\{p\notin\{u,w\}\}\cup\{u\}$.
\end{proof}

Part (ii) has a reading worth isolating, because the entire finitization
rests on it: \emph{at every time of the pair-sum lattice, the full
$n$-runner minimum and the $(n{-}1)$-runner reduced minimum coincide}.
Deleting runner $w$ costs nothing precisely at the times when $w$ is
mirrored by $u$. This is the sense in which the pair-sum lattice is the
natural finite skeleton of the continuous problem.

\subsection{Attainment at a pair-sum time}

\begin{lemma}[attainment]\label{lem:attainment}
There exist a pair $\{u,w\}$ and an integer $k^\ast$ with
$1\le k^\ast\le N-1$ such that $g(k^\ast/N)=M(V)$. Specifically, let
$t^\ast$ be any global argmax, $m:=M(V)=g(t^\ast)$, and let
$B=\{p: f_p(t^\ast)=m\}$ be the binders. Then:
\begin{enumerate}[leftmargin=1.8em,itemsep=1pt]
\item[(a)] if $0<m<\tfrac12$: some binder is \emph{rising} (its residue
  $v_p t^\ast\bmod 1$ lies in $(0,\tfrac12)$) and some binder is
  \emph{falling} (residue in $(\tfrac12,1)$); choosing $u$ rising and $w$
  falling, $(v_u+v_w)t^\ast$ is an integer, so
  $t^\ast=k^\ast/N_{uw}$;
\item[(b)] if $m=\tfrac12$, the half case of
Lemma~\ref{lem:halfcase} applies: all speeds are odd, $t^\ast=\tfrac12$,
and every pair has even $N$ with $k^\ast=N/2$;
\item[(c)] $m=0$ cannot occur.
\end{enumerate}
\end{lemma}

\begin{lemma}[the half case]\label{lem:halfcase}
If $M(V)=\tfrac12$ then every speed is odd and $t^\ast=\tfrac12$
attains it; every pair sum is even, so $t^\ast=(N/2)/N$ is a
pair-lattice time for any pair with $N\ge 4$ (such a pair exists for
$n\ge2$ unless all speeds equal $1$; in that degenerate case
$M=\tfrac12=t^\ast$ with $N=2$, $k^\ast=1$).
\end{lemma}

\begin{proof}
(c) was noted above ($M(V)>0$).

(a) Assume $0<m<\tfrac12$. Each binder $p\in B$ has residue exactly $m$
(\emph{rising}: on the arc $(0,\tfrac12)$ the position $x=v_p t\bmod 1$
satisfies $f_p=x$, which is increasing at slope $+v_p$) or exactly $1-m$
(\emph{falling}: on $(\tfrac12,1)$, $f_p=1-x$ with slope $-v_p$). No
binder sits at a peak ($f=\tfrac12$) or a valley ($f=0$) because $m$ is
neither $\tfrac12$ nor $0$.

\emph{Claim: not all binders rise, and not all binders fall.} Suppose all
binders rise; the falling case is symmetric. For non-binders let
$\delta_p=f_p(t^\ast)-m>0$. All $f_p$ are $v_p$-Lipschitz. Choose
$s>0$ with
\[
s<\min\Bigl(\;\min_{p\notin B}\delta_p/v_p,\;\;
                \min_{p\in B}\tfrac{1}{2}(1/2-m)/v_p\Bigr).
\]
Two remarks make the estimates airtight. First, the non-binder bound
uses only the \emph{global} Lipschitz property
$f_p(t\pm s)\ge f_p(t)-v_p|s|$, which holds across the cusps of
$f_p=\wnc{v_pt}$---a non-binder crossing a cusp during the
perturbation is immaterial, because the one-sided Lipschitz lower
bound is cusp-free. Second, the binder linearity window: a rising
binder's position $x_p=v_pt^\ast\bmod 1$ equals $m\in(0,\tfrac12)$,
and $f_p$ is linear on $(0,\tfrac12)$; the window survives the shift
iff $x_p+v_ps'<\tfrac12$, which the $s$-bound guarantees per binder
($v_js'\le\tfrac12\,(\tfrac12-m)<\tfrac12-m$)---this is also where
binders with \emph{different} slopes are handled: each binder gets its
own window from the same minimum.
For binders (all rising), $0<s'\le s$ keeps the position inside
$(0,\tfrac12)$ where $f_p$ is linear, so
$f_p(t^\ast+s')=m+v_p s'>m$. For non-binders,
$f_p(t^\ast+s')\ge f_p(t^\ast)-v_p s'=m+(\delta_p-v_p s')>m$ by the
choice of $s$. Hence $g(t^\ast+s')>m=M(V)$, contradicting maximality. If
all binders fall, perturb backwards ($s'<0$); the symmetric argument gives
the same contradiction. This proves the claim.

Take $u\in B$ rising and $w\in B$ falling. Then
$v_u t^\ast\equiv m$ and $v_w t^\ast\equiv -m \pmod 1$, so
$(v_u+v_w)t^\ast\equiv 0 \pmod 1$: the integer
$k^\ast:=(v_u+v_w)t^\ast$ satisfies $t^\ast=k^\ast/N_{uw}$. Since
$g(0)=0<m$, $t^\ast\ne 0$; representing $t^\ast\in(0,1)$ gives
$0<k^\ast<N$, i.e.\ $1\le k^\ast\le N-1$.

\end{proof}

\begin{proof}[Proof of Lemma~\ref{lem:halfcase}]
If $m=\tfrac12$ then \emph{every} runner binds (no $f_p$ exceeds
$\tfrac12$), i.e.\ $v_p t^\ast\equiv \tfrac12 \pmod 1$ for all $p$. Write
$t^\ast=a/b$ in lowest terms. Then $b\mid v_p-v_1$ for all $p$ (as
$\gcd(a,b)=1$), and $2v_1 a=b(2j+1)$ for some integer $j$, so $b\mid
2v_1$. If $b$ were odd, $b\mid v_1$, hence $b\mid v_p$ for all $p$, hence
$b\mid\gcd(V)=1$ and $b=1$---but then $t^\ast=0$, excluded. So $b=2c$ is
even; then $c\mid v_1$ and $v_p\equiv v_1\pmod{2c}$ forces
$v_p\equiv 0\pmod c$ for all $p$, so $c\mid\gcd(V)=1$, $c=1$, $b=2$:
$t^\ast=\tfrac12$. Then $v_1/2\equiv\tfrac12\pmod 1$ forces $v_1$ odd,
and $v_p\equiv v_1\pmod 2$ makes every speed odd. Conversely, if all
speeds are odd then $t=\tfrac12$ has $f_p=\tfrac12$ for every $p$,
which proves attainment.
\end{proof}

\begin{remark}[which speed convention is proved where]
The equality theorem and both lemmas above are proved for possibly
non-distinct speeds; nothing in this section uses distinctness. The
conjecture itself is stated for distinct speeds (WLOG after collapsing
equal speeds, which can only decrease $M$), and the covering framework
downstream assumes distinct speeds where it says so. The half case
above is the one place where non-distinctness changes the picture
(all speeds equal $1$), and it is handled inside
Lemma~\ref{lem:halfcase}.
\end{remark}

\begin{remark}[which speed convention is proved where]
The equality theorem and both lemmas above are proved for possibly
non-distinct speeds; nothing in this section uses distinctness. The
conjecture itself is stated for distinct speeds (WLOG after collapsing
equal speeds, which can only decrease $M$), and the covering framework
downstream assumes distinct speeds where it says so. The half case
above is the one place where non-distinctness changes the picture
(all speeds equal $1$), and it is handled inside
Lemma~\ref{lem:halfcase}.
\end{remark}

\begin{remark}[what forced two-sidedness means]\label{rem:twosided}
The perturbation argument is the rigorous form of an empirical
regularity: in exhaustive scans at $n=4,5,6$, every exact argmax
breakpoint has binders on both sides ($480/480$ random sets
\cite{prog-artifacts}). It also shows why the folklore's
\emph{difference} denominators $|v_u-v_w|$ never have to carry the
optimum: same-sided binder pairs are exactly the configurations that
produce difference relations, and Lemma~\ref{lem:attainment}(a) rules out
their being the only binders. The sharpening is confirmed structurally in
\S\ref{sec:equality-verify}: in $885$ test sets and, under a from-scratch
reimplementation, in $2{,}870$ further sets spanning $n=2$ to $6$ and
speeds to $200$, no unrestricted argmax lies off the pair-sum skeleton.
\end{remark}

\subsection{Proof of the theorem}

\begin{proof}[Proof of Theorem~\ref{thm:equality}]
\emph{($\ge$, the constructive bound).} Fix any pair $\{u,w\}$ and any
$1\le k\le N-1$. Lemma~\ref{lem:mirror}(ii) gives
$\min_{p\ne w}\lVert v_p k/N\rVert=g(k/N)\le M(V)$; take the max over
$k$ to get $\tau(u,w)\le M(V)$, hence
$\max_{\{u,w\}}\tau\le M(V)$.

\emph{($\le$, attainment).} By Lemma~\ref{lem:attainment} there is a pair
$\{u,w\}$ with $g(k^\ast/N)=M(V)$. Lemma~\ref{lem:mirror}(ii) again:
\[
\min_{p\ne w}\Bigl\lVert \frac{v_p k^\ast}{N}\Bigr\rVert
=g(k^\ast/N)=M(V),
\]
so $\tau(u,w)\ge M(V)$. In case (b)---all speeds odd---this works for
every pair at $k^\ast=N/2$: all runners sit at $\tfrac12$, the mirror is
trivially exact, and the reduced minimum equals $M(V)$.

Combining both directions, $\max_{\{u,w\}}\tau(u,w)=M(V)$, with equality
attained at the pair produced by Lemma~\ref{lem:attainment}. The
two-sided conclusion in the $M<\tfrac12$ case is part (a) of the lemma:
the rising binder $u$ has residue $m^\ast$ and the falling binder $w$ has
residue $-m^\ast$.
\end{proof}

\subsection{Verification and novelty status}\label{sec:equality-verify}

The theorem was verified by two deliberately non-circular computations,
recorded in the artifact bundle \cite{prog-artifacts}. The first, the
script \texttt{V1\_equality\_fresh\_check.py}, computes $M$
\emph{independently} from the folklore breakpoint set
\[
D=\{2v_p\}\cup\{v_p+v_q\}\cup\{|v_p-v_q|\},
\]
which is exhaustive for the vertices of the piecewise-linear lower
envelope and carries no pair-sum assumption, and compares it against the
pair-lattice maxima: on $885$ speed sets ($479$ at $n\le 4$ with
$v\le 12$; $393$ random at $n=5$ with $v\le 20$; $13$ targeted edge
cases including all-odd, tight, and the exceptional set
$(2,6,8,10,11)$), equality holds exactly ($885/885$), every unrestricted
argmax lies on a pair-sum lattice ($885/885$, and \emph{zero} sets have
a difference-only argmax), and the all-odd pole case is exact. The
second, \texttt{EQ\_independent\_verify.py}, is a from-scratch
reimplementation written against the theorem statement alone---a
different code path, exact integer arithmetic, time domain: it maximizes
$g$ over the full candidate set of individual breakpoints $j/(2v_p)$
together with \emph{all} pairwise crossings $c/(v_p\pm v_q)$, differences
included, and compares against the pair-sum lattices. On $2{,}870$ sets
($300$ to $1{,}200$ random primitive sets at each $n=2,\dots,6$, of which
$1{,}200$ lie in the paper's measurement domain $v\le 56$, plus a
beyond-corpus probe at $v\le 200$), equality holds exactly
($2{,}870/2{,}870$), the LRC bound holds on every set, and the anchors
are exact: the regular sets $(1,\dots,n)$ at $1/(n{+}1)$, the tight sets
$(1,2,3,5)$, $(1,3,4,7)$, $(1,3,4,5,9)$, and the base rung's open set
$(3,5,8,13)$, whose exact loneliness is $M=5/21$. The fresh scan adds
one honest nuance: on $962$ of the $2{,}870$ sets a \emph{difference}
time also attains the optimum (e.g.\ $t=\tfrac12$ for $V=(1,59)$).
Difference denominators sometimes carry an argmax, but never a needed
one; the sharpening, as Lemma~\ref{lem:attainment} states it, is that
sum lattices alone always suffice, and that is exactly what the
equality check verifies. Cross-checks with a third method (direct grid
maximization) and at $n=5,6$ on targeted tight sets bring the total to
over $3{,}700$ sets checked by at least two independent routes.

On novelty: two literature rounds ($56$ queries, the arXiv ring
enumerated for 2024--2026) found the ``sum or difference'' denominator
observation in a comment on Tao's 2015 post \cite{Tao-blog} and nothing
else at the level claimed here---but the spectrum line computes exact
optima on infinite families \cite{Cordella}, so pieces may exist in
another language. The honest position, maintained throughout: the
technique is folklore; the sum-only sharpening, the exact equality, the
$(\mathrm{TAU}$-$_n)$ statement, and everything downstream of them are
not located in the searched literature (two rounds, $56$ queries, the
arXiv ring enumerated for 2024--2026); the correct prior, given how
actively the spectrum line computes exact optima \cite{Cordella}, is
that pieces may exist in another language.
